\documentclass{article}
\usepackage{fontspec}
\setmainfont{Libertinus Serif}
\usepackage[a4paper, margin=1in]{geometry}
\usepackage{amsmath, amsthm, amssymb, mathtools, enumitem, hyperref}
\usepackage[portuguese]{babel}

\DeclareMathOperator{\im}{im}

\begin{document}

\begin{center}
    \large \textbf{MAT 01156 \\
    ÁLGEBRA LINEAR II \\}
    \vspace{0.4cm}
    \today
\end{center}

\begin{center}
    \large \textbf{Lista 4 - Soluções}
\end{center}

\begin{enumerate}
    \item [\textbf{E1}]
    \item [\textbf{E2}] Escreve \(v = \langle v, u_1 \rangle u_1 + \cdots + \langle v, u_n \rangle u_m\) e \(w = \langle w, u_1 \rangle u_1 + \cdots + \langle w, u_m \rangle u_m\). Então,
    \[
        \langle v, w \rangle = \left\langle \sum_{i=1}^m \langle v, u_i \rangle u_i, \sum_{j=1}^m \langle w, u_j \rangle u_j \right\rangle = \sum_{i=1}^m \sum_{j=1}^m \langle v, u_i \rangle \overline{\langle w, u_j \rangle} \langle u_i, u_j \rangle.
    \] Como \(\{u_1, \ldots, u_n\}\) é ortonormal, \(\langle u_i, u_j \rangle = 0\) se \(i \neq j\) e \(\langle u_i, u_i \rangle = 1\). Portanto,
    \[
        \langle v, w \rangle = \sum_{i=1}^n \langle v, u_i \rangle \overline{\langle w, u_i \rangle}.
    \]
    \item [\textbf{E3}]
    \item [\textbf{E4}]
    \item [\textbf{E5}]
    \item [\textbf{E6}]
    \item [\textbf{E7}]
    \item [\textbf{E8}]
    \item [\textbf{E9}]
    \item [\textbf{E10}]
    \item [\textbf{E11}]
\end{enumerate}

\end{document}